\section{Related Work}
\subsection{Reward Engineering for Reinforcement Learning}
Reward shaping augments task objectives with intermediate signals that improve credit assignment while preserving the optimal policy only under restricted transformations~\cite{ng1999policy}. Inverse reward design treats a specified reward as imperfect evidence of designer intent and highlights the gap between written objectives and desired behavior~\cite{hadfield2017inverse}. More broadly, reward misspecification can produce unsafe or unintended behavior even when the reward appears semantically reasonable~\cite{amodei2016concrete}. These results motivate methods that treat reward engineering as an iterative debugging problem rather than a one-shot specification step.

\subsection{LLM-Based Reward Generation and Refinement}
LLMs can map task descriptions to executable reward code. Language to Rewards demonstrates language-conditioned reward generation for robotic skill synthesis~\cite{language2rewards}. Text2Reward frames reward shaping as interpretable program synthesis from language instructions~\cite{text2reward}. Eureka combines reward-code generation with evolutionary optimization and component-level feedback~\cite{eureka}. CARD incorporates dynamic feedback into LLM-driven reward design~\cite{card}. REvolve studies reward evolution with human feedback in the loop~\cite{revolve}. A chain-of-thought reward-engineering approach shows that explicit reasoning traces can improve reward-design transparency~\cite{cotreward}. Diagnostic-driven refinement treats reward failure as a debugging problem and uses training diagnostics for targeted revision~\cite{wang2026diagnostic}. Collectively, these methods show that training feedback supports reward revision beyond one-shot generation.

\subsection{LLM Agents and Agentic Reward Engineering}
Language-agent research organizes LLMs to maintain state and act in external processes. ReAct couples reasoning with tool-mediated actions~\cite{react}. Reflexion stores verbal feedback as episodic experience for later decisions~\cite{reflexion}. Self-Refine shows how iterative self-feedback can improve generated outputs without parameter updates~\cite{selfrefine}. These mechanisms are entering reward design: RF-Agent formulates reward construction as sequential decision making~\cite{rfagent}, and STRIDE applies agentic reward engineering to humanoid locomotion~\cite{stride}. CREATE follows this direction but centers agency on instrumented policy training: the reward-engineering agent calls tools over training records, reasons over the returned evidence, and commits to a semantically localized action measured by the next complete training run.
